%!TEX program = xelatex
%!TEX encoding = ISO-8859-1
%!TEX program = xelatex
% !TEX encoding = ISO-8859-1
\def\ano{2019}
\def\semestre{2}
\def\disciplina{\'Algebra 1}
\def\turma{C}

\documentclass[12pt]{exam}

\usepackage{caption}
\usepackage{amssymb}
\usepackage{amsmath,amsfonts,amsthm,amstext}
\usepackage[brazil]{babel}
% \usepackage[latin1]{inputenc}
\usepackage{graphicx}
\graphicspath{{/home/jfreitas/GitHub_Repos/Algebra-1/Pictures/}{D:/OneDrive - unb.br/imagens-latex/}}
\usepackage{enumitem}
\usepackage{multicol}
\usepackage{answers}
\usepackage{tikz,ifthen}
\usetikzlibrary{lindenmayersystems}
\usetikzlibrary[shadings]
\Newassociation{solucao}{Solution}{ans}
\newtheorem{exercicio}{}

\setlength{\topmargin}{-1.0in}
\setlength{\oddsidemargin}{0in}
\setlength{\textheight}{10.1in}
\setlength{\textwidth}{6.5in}
\setlength{\baselineskip}{12mm}

\extraheadheight{0.7in}
\firstpageheadrule
\runningheadrule
\lhead{
        \begin{minipage}[c]{1.7cm}
        \includegraphics[width=1.7cm]{unb.pdf}
        \end{minipage}%
        \hspace{0pt}
        \begin{minipage}[c]{4in}
          {Universidade de Bras{\'\i}lia} --
          {Departamento de Matem{\'a}tica}
\end{minipage}
\vspace*{-0.8cm}
}
% \chead{Universidade de Bras{\'\i}lia - Departamento de Matem{\'a}tica}
% \rhead{}
% \vspace*{-2cm}

\extrafootheight{.5in}
\footrule
\lfoot{\disciplina\ - \semestre$^o$/\ano}
\cfoot{}
\rfoot{P\'agina \thepage\ de \numpages}

\newcounter{exercicios}
\renewcommand{\theexercicios}{\arabic{exercicios}}

\newenvironment{questao}[1]{
\refstepcounter{exercicios}
\ifx&#1&
\else
   \label{#1}
\fi
\noindent\textbf{Exerc{\'\i}cio {\theexercicios}:}
}

\newcommand{\resp}[1]{
\noindent{\bf Exerc{\'\i}cio #1: }}
\newcommand{\nsub}{\varsubsetneq}
\newcommand{\vaz}{\emptyset}
\newcommand{\im}{{\rm Im\,}}
\newcommand{\sub}{\subseteq}
\newcommand{\n}{\mathbb{N}}
\newcommand{\z}{\mathbb{Z}}
\newcommand{\rac}{\mathbb{Q}}
\newcommand{\real}{\mathbb{R}}
\newcommand{\complex}{\mathbb{C}}
\newcommand{\cp}[1]{\mathbb{#1}}
\newcommand{\ch}{\mbox{\textrm{car\,}}\nobreak}
\newcommand{\vesp}[1]{\vspace{ #1  cm}}
\newcommand{\compcent}[1]{\vcenter{\hbox{$#1\circ$}}}
\newcommand{\comp}{\mathbin{\mathchoice
{\compcent\scriptstyle}{\compcent\scriptstyle}
{\compcent\scriptscriptstyle}{\compcent\scriptscriptstyle}}}

\begin{document}

\begin{center}
{\Large\bf \disciplina\ - Turma \turma\ -- \semestre$^{o}$/\ano} \\ \vspace{9pt} {\large\bf
$1^{\underline{o}}$ Teste - M\'odulo 2 - Resolu\c{c}\~ao}\\
\vspace{9pt} Prof. Jos{\'e} Ant{\^o}nio O. Freitas
\end{center}
\hrule

\vspace{.6cm}

Seja $f : \real - \left\{-\dfrac{d}{c}\right\} \to \real - \left\{\dfrac{a}{c}\right\}$ dada por
\[
	f(x) =  \dfrac{ax + b}{cx + d},
\]
onde $a$, $b$, $c$, $d$ s{\~a}o n{\'u}meros reais constantes, $ad - bc \ne 0$.
\vspace{.5cm}

\questao Mostre que $f$ \'e injetora.

\noindent\textbf{Solu\c{c}\~ao:} Sejam $x_1$, $x_2 \in \real - \left\{-\dfrac{d}{c}\right\}$ tais que $f(x_1) = f(x_2)$. Temos
\begin{align*}
	f(x_1) &= f(x_2)\\
	\dfrac{ax_1 + b}{cx_1 + d} &= \dfrac{ax_2 + b}{cx_2 + d}\\
	(ax_1 + b)(cx_2 + d) &= (ax_2 + b)(cx_1 + d)\\
	acx_1x_2 + adx_1 + bcx_2 + bd &= acx_1x_2 + adx_2 + bcx_1 + bd\\
	adx_1 + bcx_2 - adx_2 - bcx_1 &= 0\\
	ad(x_1 - x_2) - bc(x_1 - x_2)&= 0\\
	(ad - bc)(x_1 - x_2) &= 0
\end{align*}
Como $ad - bc \ne 0$, ent\~ao $x_1 - x_2 = 0$. Logo $x_1 = x_2$. Portanto $f$ \'e injetora.

\vspace{.5cm}

\questao Mostre que $f$ \'e sobrejetora.

\noindent\textbf{Solu\c{c}\~ao:} A fun\c{c}\~ao $f$ ser\'a sobrejetora se para todo $y \in \real - \left\{\dfrac{a}{c}\right\}$ existir $x \in \real - \left\{-\dfrac{d}{c}\right\}$ tal que $f(x) = y$. Assim precisamos determinar se a equa\c{c}\~ao $f(x) = y$ tem solu\c{c}\~ao para todo $y \in \real - \left\{\dfrac{a}{c}\right\}$.
Temos
\begin{align*}
	f(x) &= y\\
	\dfrac{ax + b}{cx + d} &= y\\
	ax + b &= y(cx + d)\\
	ax + b &= ycx + dy\\
	x(a - yc) &= dy - b
\end{align*}
Mas na \'ultima equa\c{c}\~ao como $y \in \real - \left\{\dfrac{a}{c}\right\}$ ent\~ao $a - yc \ne 0$ logo podemos isolar $x$ obtendo
\[
	x = \dfrac{dy - b}{a - yc} \in \real - \left\{-\dfrac{d}{c}\right\}.
\]

Assim dado $y \in  \real - \left\{\dfrac{a}{c}\right\}$ tome $x = \dfrac{dy - b}{a - yc} \in \real - \left\{-\dfrac{d}{c}\right\}$. Da{\'\i}
\begin{align*}
	f(x) &= f\left(\dfrac{dy - b}{a - yc}\right) \\
	&= \dfrac{a\left(\dfrac{dy - b}{a - yc}\right) + b}{c\left(\dfrac{dy - b}{a - yc}\right) + d} \\
	&= \dfrac{\dfrac{ady - ab}{a - yc} + b}{\dfrac{cdy - cb}{a - yc} + d} \\
	&= \dfrac{\dfrac{ady - ab + ba - byc}{a - yc}}{\dfrac{cdy - cb + da - dyc}{a - yc}} \\
	&= \dfrac{(ad - bc)y}{ad - bc} \\ 
	&= y
\end{align*}
Portanto $f$ \'e sobrejetora.

\end{document}